\section{Transcendence}
\label{section-transcendence}

\noindent
We recall the standard definitions.

\begin{definition}
\label{definition-transcendence}
Let $K/k$ be a field extension.
\begin{enumerate}
\item A collection of elements $\{x_i\}_{i \in I}$ of $K$ is called
{\it algebraically independent} over $k$ if the map
$$
k[X_i; i\in I] \longrightarrow K
$$
which maps $X_i$ to $x_i$ is injective.
\item The field of fractions of a polynomial ring
$k[x_i; i \in I]$ is denoted $k(x_i; i\in I)$.
\item A {\it purely transcendental extension} of $k$ is any
field extension $K/k$ isomorphic to the field of
fractions of a polynomial ring over $k$.
\item A {\it transcendence basis} of $K/k$ is a
collection of elements $\{x_i\}_{i \in I}$ which are
algebraically independent over $k$ and such that
the extension $K/k(x_i; i\in I)$ is algebraic.
\end{enumerate}
\end{definition}

\begin{example}
\label{example-pi-transcendental}
The field $\mathbf{Q}(\pi)$ is purely transcendental because
$\pi$ isn't the root of a nonzero polynomial with rational coefficients.
In particular, $\mathbf{Q}(\pi) \cong \mathbf{Q}(x)$.
\end{example}

\begin{lemma}
\label{lemma-transcendence-degree}
Let $E/F$ be a field extension. A transcendence basis of $E$ over $F$ exists.
Any two transcendence bases have the same cardinality.
\end{lemma}

\begin{proof}
Let $A$ be an algebraically independent subset of $E$. Let $G$ be a subset
of $E$ containing $A$ that generates $E/F$. We claim we can find a
transcendence basis $B$ such that $A \subset B \subset G$.
To prove this, consider the collection $\mathcal{B}$ of algebraically
independent subsets whose members are subsets of $G$ that contain $A$.
Define a partial ordering on $\mathcal{B}$ using inclusion.
Then $\mathcal{B}$ contains at least one element $A$.
The union of the elements of a totally ordered subset $T$ of $\mathcal{B}$
is an algebraically independent subset of $E$ over $F$ since any algebraic
dependence relation would have occurred in one of the elements of $T$
(since polynomials only involve finitely many variables). The union also
contains $A$ and is contained in $G$. By Zorn's lemma, there is a maximal
element $B \in \mathcal{B}$. Now we claim $E$ is algebraic over $F(B)$.
This is because if it wasn't then there would be an element
$f \in G$ transcendental over $F(B)$ since $F(G) = E$. Then
$B \cup\{f\}$ would be algebraically independent contradicting the
maximality of $B$. Thus $B$ is our transcendence basis.

\medskip\noindent
Let $B$ and $B'$ be two transcendence bases. Without loss of generality, we
can assume that $|B'| \leq |B|$. Now we divide the proof into two cases: the
first case is that $B$ is an infinite set. Then for each $\alpha \in B'$,
there is a finite set $B_{\alpha} \subset B$ such that $\alpha$ is algebraic
over
$F(B_{\alpha})$ since any algebraic dependence relation only uses finitely many
indeterminates. Then we define $B^* = \bigcup_{\alpha\in B'} B_{\alpha}$.
By construction, $B^* \subset B$, but we claim that in fact the two sets are
equal. To see this, suppose that they are not equal, say there is an element
$\beta \in B \setminus B^*$. We know $\beta$ is algebraic over $F(B')$ which
is algebraic over $F(B^*)$. Therefore $\beta$ is algebraic over $F(B^*)$, a
contradiction. So $|B| \leq |\bigcup_{\alpha \in B'} B_{\alpha}|$.
Now if $B'$ is finite, then so is $B$ so we can assume $B'$ is infinite;
this means
$$
|B| \leq |\bigcup\nolimits_{\alpha \in B'} B_{\alpha}| = |B'|
$$
because each $B_\alpha$ is finite and $B'$ is infinite. Therefore in the
infinite case, $|B| = |B'|$.

\medskip\noindent
Now we need to look at the case where $B$ is finite.
In this case, $B'$ is also finite, so suppose
$B = \{\alpha_1, \ldots, \alpha_n\}$ and
$B' = \{\beta_1, \ldots, \beta_m\}$ with $m \leq n$.
We perform induction on $m$: if $m = 0$ then $E/F$ is algebraic so
$B = \emptyset$ so $n = 0$. If $m > 0$, there is an irreducible polynomial
$f \in F[x, y_1, \ldots, y_n]$ such that
$f(\beta_1, \alpha_1, \ldots, \alpha_n) = 0$ and such that $x$ occurs in $f$.
Since $\beta_1$ is not algebraic over $F$, $f$ must involve some $y_i$
so without loss of generality, assume $f$ uses $y_1$.
Let $B^* = \{\beta_1, \alpha_2, \ldots, \alpha_n\}$.
We claim that $B^*$ is a basis for $E/F$. To prove this claim, we see that
we have a tower of algebraic extensions
$$
E/ F(B^*, \alpha_1) / F(B^*)
$$
since $\alpha_1$ is algebraic over $F(B^*)$.
Now we claim that $B^*$ (counting multiplicity of elements) is
algebraically independent over $F$ because if it weren't, then there would be an
irreducible $g\in F[x, y_2, \ldots, y_n]$ such that
$g(\beta_1, \alpha_2, \ldots, \alpha_n) = 0$
which must involve $x$ making $\beta_1$
algebraic over $F(\alpha_2, \ldots, \alpha_n)$ which would make $\alpha_1$
algebraic over $F(\alpha_2, \ldots, \alpha_n)$ which is impossible.
So this means that $\{\alpha_2, \ldots, \alpha_n\}$ and
$\{\beta_2, \ldots, \beta_m\}$ are bases for $E$ over $F(\beta_1)$
which means by induction, $m = n$.
\end{proof}

\begin{definition}
\label{definition-transcendence-degree}
Let $K/k$ be a field extension.
The {\it transcendence degree} of $K$ over $k$ is
the cardinality of a transcendence basis of $K$ over $k$.
It is denoted $\text{trdeg}_k(K)$.
\end{definition}

\begin{lemma}
\label{lemma-transcendence-degree-tower}
Let $L/K/k$ be field extensions.
Then
$$
\text{trdeg}_k(L) =
\text{trdeg}_K(L) +
\text{trdeg}_k(K).
$$
\end{lemma}

\begin{proof}
Choose a transcendence basis $A \subset K$ of $K$ over $k$.
Choose a transcendence basis $B \subset L$ of $L$ over $K$.
Then it is straightforward to see that $A \cup B$ is a transcendence
basis of $L$ over $k$.
\end{proof}

\begin{example}
\label{example-pi-e-transcendental}
Consider the field extension $\mathbf{Q}(e, \pi)$ formed by
adjoining the numbers $e$ and $\pi$. This field extension has transcendence
degree at least $1$ since both $e$ and $\pi$ are transcendental over the
rationals. However, this field extension might have transcendence
degree $2$ if $e$ and $\pi$ are algebraically independent. Whether or
not this is true is unknown and whence the problem of determining
$\text{trdeg}(\mathbf{Q}(e, \pi))$ is open.
\end{example}

\begin{example}
\label{example-function-field-not-unique-transcendence-basis}
Let $F$ be a field and $E = F(t)$. Then $\{t\}$ is a
transcendence basis since $E = F(t)$. However, $\{t^2\}$
is also a transcendence basis since $F(t)/F(t^2)$ is algebraic.
This illustrates that while we can always decompose an extension
$E/F$ into an algebraic extension $E/F'$ and a
purely transcendental extension $F'/F$, this decomposition is not unique and
depends on choice of transcendence basis.
\end{example}

\begin{example}
\label{example-riemann-surface-transcendence}
Let $X$ be a compact Riemann surface. Then the function field
$\mathbf{C}(X)$ (see Example \ref{example-field-of-meromorphic-functions})
has transcendence degree one over $\mathbf{C}$. In fact, {\it any}
finitely generated extension of $\mathbf{C}$ of transcendence degree
one arises from a Riemann surface. There is even an equivalence of
categories between the category of compact Riemann surfaces and
(non-constant) holomorphic maps and the opposite of the category of finitely
generated extensions of $\mathbf{C}$ of transcendence degree $1$
and morphisms of $\mathbf{C}$-algebras. See \cite{Forster}.

\medskip\noindent
There is an algebraic version of the above statement as well. Given an
(irreducible) algebraic curve in projective space over an algebraically
closed field $k$ (e.g. the complex numbers), one can consider its
``field of rational functions'': basically, functions that look like
quotients of polynomials, where the denominator does not identically
vanish on the curve. There is a similar anti-equivalence of categories
(Algebraic Curves, Theorem \ref{curves-theorem-curves-rational-maps})
between smooth projective curves and
non-constant morphisms of curves and finitely generated extensions of $k$ of
transcendence degree one. See \cite{H}.
\end{example}

\begin{definition}
\label{definition-algebraically-closed-in}
Let $K/k$ be a field extension.
\begin{enumerate}
\item The {\it algebraic closure of $k$ in $K$} is the subfield
$k'$ of $K$ consisting of elements of $K$ which are algebraic over $k$.
\item We say $k$ is {\it algebraically closed in $K$} if
every element of $K$ which is algebraic over $k$ is
contained in $k$.
\end{enumerate}
\end{definition}

\begin{lemma}
\label{lemma-purely-transcendental-degree}
Let $k'/k$ be a finite extension of fields. Let
$k'(x_1, \ldots, x_r)/k(x_1, \ldots, x_r)$ be the
induced extension of purely transcendental extensions.
Then $[k'(x_1, \ldots, x_r) : k(x_1, \ldots, x_r)] = [k' : k] < \infty$.
\end{lemma}

\begin{proof}
By multiplicativity of degrees of extensions
(Lemma \ref{lemma-multiplicativity-degrees})
it suffices to prove this when $k'$ is generated by a single element
$\alpha \in k'$ over $k$. Let $f \in k[T]$ be the minimal polynomial
of $\alpha$ over $k$. Then $k'(x_1, \ldots, x_r)$ is generated
by $\alpha, x_1, \ldots, x_r$ over $k$ and hence $k'(x_1, \ldots, x_r)$
is generated by $\alpha$ over $k(x_1, \ldots, x_r)$.
Thus it suffices to show that $f$ is still irreducible as
an element of $k(x_1, \ldots, x_r)[T]$. We only sketch the proof.
It is clear that $f$ is irreducible as an element of
$k[x_1, \ldots, x_r, T]$ for example because $f$ is monic
as a polynomial in $T$ and any putative factorization in
$k[x_1, \ldots, x_r, T]$ would lead to a factorization in $k[T]$
by setting $x_i$ equal to $0$. By Gauss' lemma we conclude.
\end{proof}

\begin{lemma}
\label{lemma-algebraic-closure-in-finitely-generated}
Let $K/k$ be a finitely generated field extension.
The algebraic closure of $k$ in $K$ is finite over $k$.
\end{lemma}

\begin{proof}
Let $x_1, \ldots, x_r \in K$ be a transcendence basis for $K$
over $k$. Then $n = [K : k(x_1, \ldots, x_r)] < \infty$.
Suppose that $k \subset k' \subset K$ with $k'/k$ finite.
In this case
$[k'(x_1, \ldots, x_r) : k(x_1, \ldots, x_r)] = [k' : k] < \infty$, see
Lemma \ref{lemma-purely-transcendental-degree}.
Hence
$$
[k' : k] = [k'(x_1, \ldots, x_r) : k(x_1, \ldots, x_r)] \leq
[K : k(x_1, \ldots, x_r)] = n.
$$
In other words, the degrees of finite subextensions are bounded
and the lemma follows.
\end{proof}






